\section{Conclusão}

% ---------------------------------------------------------------- 35
\begin{frame}{Redes neurais e operadores}
  Em resumo, o ganhador é quando juntamos física com dados.
  \textbf{Redes neurais (MLP e PINN)}
  \begin{itemize}[<+->]
    \item \textbf{MLP, dados puros.} Reproduz a onda; banda baixa depressa, alta
          devagar;
    \item \textbf{PINN, dados e física.} Menor erro relativo entre os três;
    \item \textbf{PINN, física pura.} Sem dados, não aprende detalhes da onda.
  \end{itemize}
  \vspace{2pt}
  \textbf{Operadores (FNO e PINO)}
  \begin{itemize}[<+->]
    \item \textbf{FNO.} Ajusta bem os dados, mas com viés espectral;
    \item \textbf{PINO, física pura.} Escapam as ondulações de alta frequência;
    \item \textbf{PINO, dados e física.} O caso mais consistente na varredura;
    \item \textbf{Invariância à discretização.} Fora da grade de treino, queda modesta.
  \end{itemize}
\end{frame}

% ---------------------------------------------------------------- 37
\begin{frame}{Limitações}
  Temos que levar em consideração que:
  \begin{itemize}[<+->]
    \item Utilizamos uma única condição inicial;
    \item Consideramos $\alpha = \beta$;
    \item A solução de referência é numérica;
    \item O estudo do viés espectral é bastante simplificado.
    
    \uncover<+->{
      \begin{information}[Viés espectral]
        O estudo de viés espectral que fizemos é um cenário muito simplificado, não pode-se garantir que as estimativas encontradas se estendam para outras arquiteturas.
      \end{information}
    }
  \end{itemize}
\end{frame}

% ---------------------------------------------------------------- 38
\begin{frame}{Considerações finais}
  Vimos que o viés espectral é um comportamento real dessas arquiteturas, quem as implementa deve estar, ao menos, ciente disso.
  \begin{itemize}[<+->]
    \item Nada neste trabalho favorece que ML substitui métodos clássicos de solução de EDPs;
    \item Métodos clássicos vêm com a vantagem teórica de estimativas de erros e garantias de convergência;
    \item As comparações deste trabalho são amarradas apenas para o sistema de Boussinesq, não se pode generalizar para outras EDPs.
  \end{itemize}
  \begin{information}[Fechamento]
    Conforme grande crescimento de ML, esperamos que este trabalho sirva como guia para estudantes que desejam conhecer a aplicação de SciML à EDPs.
  \end{information}
\end{frame}
